\documentclass[10pt,a4paper]{amsart}
\usepackage{amsmath,amssymb,amsthm,graphicx}
\usepackage{verbatim}
\newenvironment{reference}{\comment}{\endcomment}
%\newenvironment{reference}{}{}
\newenvironment{slogan}{\comment}{\endcomment}
\newenvironment{history}{\comment}{\endcomment}

% For commutative diagrams we use Xy-pic
\usepackage[all]{xy}

% We use 2cell for 2-commutative diagrams.
\xyoption{2cell}
% Theorem environments.
%
\theoremstyle{plain}
\newtheorem{theorem}[subsection]{Theorem}
\newtheorem{proposition}[subsection]{Proposition}
\newtheorem{lemma}[subsection]{Lemma}

\theoremstyle{definition}
\newtheorem{definition}[subsection]{Definition}
\newtheorem{example}[subsection]{Example}
\newtheorem{exercise}[subsection]{Exercise}
\newtheorem{situation}[subsection]{Situation}

\theoremstyle{remark}
\newtheorem{remark}[subsection]{Remark}
\newtheorem{remarks}[subsection]{Remarks}

\numberwithin{equation}{subsection}

% Macros
%
\def\lim{\mathop{\mathrm{lim}}\nolimits}
\def\colim{\mathop{\mathrm{colim}}\nolimits}
\def\Spec{\mathop{\mathrm{Spec}}}
\def\Hom{\mathop{\mathrm{Hom}}\nolimits}
\def\Ext{\mathop{\mathrm{Ext}}\nolimits}
\def\SheafHom{\mathop{\mathcal{H}\!\mathit{om}}\nolimits}
\def\SheafExt{\mathop{\mathcal{E}\!\mathit{xt}}\nolimits}
\def\Sch{\mathit{Sch}}
\def\Mor{\mathop{\mathrm{Mor}}\nolimits}
\def\Ob{\mathop{\mathrm{Ob}}\nolimits}
\def\Sh{\mathop{\mathit{Sh}}\nolimits}
\def\NL{\mathop{N\!L}\nolimits}
\def\CH{\mathop{\mathrm{CH}}\nolimits}
\def\proetale{{pro\text{-}\acute{e}tale}}
\def\etale{{\acute{e}tale}}
\def\QCoh{\mathit{QCoh}}
\def\Ker{\mathop{\mathrm{Ker}}}
\def\Im{\mathop{\mathrm{Im}}}
\def\Coker{\mathop{\mathrm{Coker}}}
\def\Coim{\mathop{\mathrm{Coim}}}

% Boxtimes
%
\DeclareMathSymbol{\boxtimes}{\mathbin}{AMSa}{"02}

%
% Macros for moduli stacks/spaces
%
\def\QCohstack{\mathcal{QC}\!\mathit{oh}}
\def\Cohstack{\mathcal{C}\!\mathit{oh}}
\def\Spacesstack{\mathcal{S}\!\mathit{paces}}
\def\Quotfunctor{\mathrm{Quot}}
\def\Hilbfunctor{\mathrm{Hilb}}
\def\Curvesstack{\mathcal{C}\!\mathit{urves}}
\def\Polarizedstack{\mathcal{P}\!\mathit{olarized}}
\def\Complexesstack{\mathcal{C}\!\mathit{omplexes}}
% \Pic is the operator that assigns to X its picard group, usage \Pic(X)
% \Picardstack_{X/B} denotes the Picard stack of X over B
% \Picardfunctor_{X/B} denotes the Picard functor of X over B
\def\Pic{\mathop{\mathrm{Pic}}\nolimits}
\def\Picardstack{\mathcal{P}\!\mathit{ic}}
\def\Picardfunctor{\mathrm{Pic}}
\def\Deformationcategory{\mathcal{D}\!\mathit{ef}}
\usepackage{cleveref}
\crefname{originallemma}{original lemma}{original lemmas}
\Crefname{originallemma}{Original lemma}{Original lemmas}
\makeatletter
\newlabel{varieties-lemma-locally-Noetherian-base-change}{{038R}{source}{}{sourcepointer-1}{}}
\newlabel{varieties-lemma-locally-Noetherian-base-change@cref}{{[originallemma][0][]038R}{[1][0][]source}{}{}{}}
\makeatother
\crefname{originallemma}{original lemma}{original lemmas}
\Crefname{originallemma}{Original lemma}{Original lemmas}
\makeatletter
\newlabel{algebra-lemma-base-change-flat-up-down}{{00MQ}{source}{}{sourcepointer-2}{}}
\newlabel{algebra-lemma-base-change-flat-up-down@cref}{{[originallemma][0][]00MQ}{[1][0][]source}{}{}{}}
\makeatother
\crefname{originallemma}{original lemma}{original lemmas}
\Crefname{originallemma}{Original lemma}{Original lemmas}
\makeatletter
\newlabel{lemma-gorenstein-base-change}{{0C03}{source}{}{sourcepointer-3}{}}
\newlabel{lemma-gorenstein-base-change@cref}{{[originallemma][0][]0C03}{[1][0][]source}{}{}{}}
\makeatother
\title{Gorensteinness at a point is stable under base change and reflects it under flatness hypotheses}
\author{The Stacks Project Authors}
\date{stacks-project:duality.tex:local-snapshot-20261012; GNU FDL 1.2 or later}
\begin{document}
\maketitle
\noindent\textit{Editorial note.} All delivered mathematical text is the author's own, copied from the pinned Stacks Project chapter file duality.tex; no statement or proof step was rewritten, reordered or removed.\par\medskip
\noindent\textit{Editorial note.} Preamble provenance: the pinned chapter file opens with an input of the shared preamble and carries no preamble text of its own, so the project's single Stacks-wide preamble staged at sources/stacks-duality/preamble.tex is excerpted verbatim under the recipe fixed by LANE\_BRIEF 22.1, as three spans, lines 8-13 (verbatim package and the reference/slogan/history environments), lines 14-18 (Xy-pic setup) and lines 154-217 (the native theorem declarations and mathematical macros). Nothing was added, removed or redefined; the wrapper class and the amsmath/amssymb/amsthm/graphicx base packages are supplied by the assembler.\par\medskip
\noindent\textit{Editorial note.} The counted statement and proof contain no cross reference, no citation and no included file, so the delivered body is self contained.\par\medskip
\noindent\textit{Editorial note.} Canonical location of the counted proposition: Stacks tag 0C07 (sources/stacks-duality/pinned-original-source.tex line 6614).\par\medskip
\noindent\textit{Allowed standard prerequisites (editorial).}\begin{itemize}
\item The definition of Gorensteinness at a point in terms of flatness and the fibre. Same chapter, the definition of a Gorenstein morphism preceding this lemma, invoked in the statement. Fixes the two conditions checked in each implication.
\item Base change of the Gorenstein property of a ring. Same chapter, Lemma labelled lemma-gorenstein-base-change, cited at proof lines 6652-6656 and 6662-6666 and 6672-6676. Transfers Gorensteinness of the fibre in both directions.
\item Flatness base change up and down. Algebra, Lemma labelled algebra-lemma-base-change-flat-up-down, cited at proof lines 6650-6652 and 6668-6672. Transfers flatness of the local ring maps in all three implications.
\item Local Noetherianity of the fibres after a finite type base change. Varieties, Lemma labelled varieties-lemma-locally-Noetherian-base-change, cited at proof lines 6632-6636. Makes the statement meaningful.
\end{itemize}\medskip
\ifdefined\hypertarget\hypertarget{sourcepointer-1}{}\fi
\noindent\textit{Source reference.} Original lemma 038R: sources/stacks-chow/tags.txt, source lines 4246--4246. Cross-chapter standard prerequisite located in the project tag list staged in the local snapshot; invoked at proof line 6635 of the counted proof. It is retained only as a navigation pointer; every load-bearing step is inside the delivered span.\par\medskip
\ifdefined\hypertarget\hypertarget{sourcepointer-2}{}\fi
\noindent\textit{Source reference.} Original lemma 00MQ: sources/stacks-chow/tags.txt, source lines 832--832. Cross-chapter standard prerequisite located in the project tag list staged in the local snapshot; invoked at proof line 6651 of the counted proof. It is retained only as a navigation pointer; every load-bearing step is inside the delivered span.\par\medskip
\ifdefined\hypertarget\hypertarget{sourcepointer-3}{}\fi
\noindent\textit{Source reference.} Original lemma 0C03: sources/stacks-duality/pinned-original-source.tex, source lines 6480--6480. Same-chapter pointer into the author environment lemma at source line 6480; invoked at proof line 6656 of the counted proof. It is retained only as a navigation pointer; every load-bearing step is inside the delivered span.\par\medskip
\begin{lemma}
\label{lemma-base-change-gorenstein}
Let $f : X \to Y$ be a morphism of schemes.
Assume that all the fibres $X_y$ are locally Noetherian schemes.
Let $Y' \to Y$ be locally of finite type. Let $f' : X' \to Y'$
be the base change of $f$.
Let $x' \in X'$ be a point with image $x \in X$.
\begin{enumerate}
\item If $f$ is Gorenstein at $x$, then
$f' : X' \to Y'$ is Gorenstein at $x'$.
\item If $f$ is flat at $x$ and $f'$ is Gorenstein at $x'$, then $f$
is Gorenstein at $x$.
\item If $Y' \to Y$ is flat at $f'(x')$ and $f'$ is Gorenstein at
$x'$, then $f$ is Gorenstein at $x$.
\end{enumerate}
\end{lemma}

\begin{proof}
Note that the assumption on $Y' \to Y$ implies that for $y' \in Y'$
mapping to $y \in Y$ the field extension $\kappa(y')/\kappa(y)$
is finitely generated. Hence also all the fibres
$X'_{y'} = (X_y)_{\kappa(y')}$ are locally Noetherian, see
Varieties, Lemma \ref{varieties-lemma-locally-Noetherian-base-change}.
Thus the lemma makes sense. Set $y' = f'(x')$ and $y = f(x)$.
Hence we get the following commutative diagram of local rings
$$
\xymatrix{
\mathcal{O}_{X', x'} & \mathcal{O}_{X, x} \ar[l] \\
\mathcal{O}_{Y', y'} \ar[u] & \mathcal{O}_{Y, y} \ar[l] \ar[u]
}
$$
where the upper left corner is a localization of the tensor product
of the upper right and lower left corners over the lower right corner.

\medskip\noindent
Assume $f$ is Gorenstein at $x$.
The flatness of $\mathcal{O}_{Y, y} \to \mathcal{O}_{X, x}$
implies the flatness of $\mathcal{O}_{Y', y'} \to \mathcal{O}_{X', x'}$, see
Algebra, Lemma \ref{algebra-lemma-base-change-flat-up-down}.
The fact that $\mathcal{O}_{X, x}/\mathfrak m_y\mathcal{O}_{X, x}$
is Gorenstein implies that
$\mathcal{O}_{X', x'}/\mathfrak m_{y'}\mathcal{O}_{X', x'}$
is Gorenstein, see
Lemma \ref{lemma-gorenstein-base-change}. Hence we see that $f'$
is Gorenstein at $x'$.

\medskip\noindent
Assume $f$ is flat at $x$ and $f'$ is Gorenstein at $x'$.
The fact that $\mathcal{O}_{X', x'}/\mathfrak m_{y'}\mathcal{O}_{X', x'}$
is Gorenstein implies that
$\mathcal{O}_{X, x}/\mathfrak m_y\mathcal{O}_{X, x}$
is Gorenstein, see
Lemma \ref{lemma-gorenstein-base-change}. Hence we see that $f$
is Gorenstein at $x$.

\medskip\noindent
Assume $Y' \to Y$ is flat at $y'$ and $f'$ is Gorenstein at
$x'$. The flatness of $\mathcal{O}_{Y', y'} \to \mathcal{O}_{X', x'}$
and $\mathcal{O}_{Y, y} \to \mathcal{O}_{Y', y'}$ implies the flatness
of $\mathcal{O}_{Y, y} \to \mathcal{O}_{X, x}$, see
Algebra, Lemma \ref{algebra-lemma-base-change-flat-up-down}.
The fact that $\mathcal{O}_{X', x'}/\mathfrak m_{y'}\mathcal{O}_{X', x'}$
is Gorenstein implies that
$\mathcal{O}_{X, x}/\mathfrak m_y\mathcal{O}_{X, x}$
is Gorenstein, see
Lemma \ref{lemma-gorenstein-base-change}. Hence we see that $f$
is Gorenstein at $x$.
\end{proof}

\end{document}
